% GENERATED FILE. Edit canonical lessons, then run npm run generate:books.
% canonical-source-hash: fnv1a64:e738c4ec038d821c
% canonical-lessons: ML-C445-anivaryamaya, ML-C445-nirnnayakamaya, ML-C445-atisthanaparamaya, ML-C445-gurutaramaya, ML-C445-sankirnnamaya

\chapter{How Serious}
\label{ch:ml-how-serious}
\hlchaptermodality{445}

\begin{chapteropening}
\textbf{By the end of this chapter:} \emph{I can say inevitable, crucial, fundamental, serious and complex.}

Complete the last lesson of chapter 445: I can say inevitable, crucial, fundamental, serious and complex.
\end{chapteropening}

\section[\texorpdfstring{\ml{അനിവാര്യമായ} (anivāryamāya)}{അനിവാര്യമായ (anivāryamāya)}]{\ml{അനിവാര്യമായ} (anivāryamāya) — inevitable; essential}

\label{lesson:ML-C445-anivaryamaya}



\begin{quote}
Before the new one: say the Malayalam for as per, then the Malayalam for a role.
\end{quote}

\subsection*{You'll want to know: \ml{അനിവാര്യമായ}}
\textbf{\ml{അനിവാര്യമായ}} — \emph{anivāryamāya} — \textquotedblleft{}inevitable; essential\textquotedblright{}.

\begin{quote}
\ml{ആ} \ml{മാറ്റം} \ml{അനിവാര്യമായിരുന്നു}.
\end{quote}

\begin{quote}
\emph{ā māṟṟaṁ anivāryamāyirunnu.}
\end{quote}

\begin{quote}
That change was inevitable.
\end{quote}

As a predicate: \ml{അത്} \ml{അനിവാര്യമാണ്}.

\begin{cousinweb}
\textbf{\ml{അനിവാര്യമായ}} (\emph{anivāryamāya}) is built from Sanskrit \emph{anivārya}, \textquotedblleft{}not to be warded off, unavoidable\textquotedblright{} (Monier-Williams), with the adjective ending -\ml{ആയ}.
\end{cousinweb}

\begin{grammarlens}[title={What you've built}]
One more word for saying what happens if.
\end{grammarlens}

\subsection*{Your turn}
\begin{itemize}
\raggedright
  \item \emph{Say it:} \emph{anivāryamāya}
  \item \emph{Say it:} \emph{anivāryamāya}, once more
  \item \emph{Say it:} \emph{ā māṟṟaṁ anivāryamāyirunnu.} — That change was inevitable.
  \item \emph{Recall:} say the Malayalam for as per, then the Malayalam for a role, then say \emph{anivāryamāya} again
\end{itemize}

\begin{tcolorbox}[breakable,colback=teal!4,colframe=teal!35!black,title={Before you move on}]
What does \ml{അനിവാര്യമായ} mean? (\textquotedblleft{}Inevitable; essential\textquotedblright{}.) Say it once more.
\end{tcolorbox}
\section[\texorpdfstring{\ml{നിർണ്ണായകമായ} (nirṇṇāyakamāya)}{നിർണ്ണായകമായ (nirṇṇāyakamāya)}]{\ml{നിർണ്ണായകമായ} (nirṇṇāyakamāya) — crucial; decisive}

\label{lesson:ML-C445-nirnnayakamaya}



\begin{quote}
Before the new one: say the Malayalam for a role, then the Malayalam for inevitable.
\end{quote}

\subsection*{You'll want to know: \ml{നിർണ്ണായകമായ}}
\textbf{\ml{നിർണ്ണായകമായ}} — \emph{nirṇṇāyakamāya} — \textquotedblleft{}crucial; decisive\textquotedblright{}.

\begin{quote}
\ml{നാളത്തെ} \ml{യോഗം} \ml{വളരെ} \ml{നിർണ്ണായകമാണ്}.
\end{quote}

\begin{quote}
\emph{nāḷatte yōgaṁ vaḷare nirṇṇāyakamāṇŭ.}
\end{quote}

\begin{quote}
Tomorrow's meeting is very crucial.
\end{quote}

As a predicate: \ml{നിർണ്ണായകമാണ്}.

\begin{cousinweb}
\textbf{\ml{നിർണ്ണായകമായ}} (\emph{nirṇṇāyakamāya}) is built from Sanskrit \emph{nirṇāyaka}, \textquotedblleft{}settling, conclusive\textquotedblright{} (Monier-Williams), with the adjective ending -\ml{ആയ}: what settles the matter.
\end{cousinweb}

\begin{grammarlens}[title={What you've built}]
One more word for saying what happens if.
\end{grammarlens}

\subsection*{Your turn}
\begin{itemize}
\raggedright
  \item \emph{Say it:} \emph{nirṇṇāyakamāya}
  \item \emph{Say it:} \emph{nirṇṇāyakamāya}, once more
  \item \emph{Say it:} \emph{nāḷatte yōgaṁ vaḷare nirṇṇāyakamāṇŭ.} — Tomorrow's meeting is very crucial.
  \item \emph{Recall:} say the Malayalam for a role, then the Malayalam for inevitable, then say \emph{nirṇṇāyakamāya} again
\end{itemize}

\begin{tcolorbox}[breakable,colback=teal!4,colframe=teal!35!black,title={Before you move on}]
What does \ml{നിർണ്ണായകമായ} mean? (\textquotedblleft{}Crucial; decisive\textquotedblright{}.) Say it once more.
\end{tcolorbox}
\section[\texorpdfstring{\ml{അടിസ്ഥാനപരമായ} (aṭisthānaparamāya)}{അടിസ്ഥാനപരമായ (aṭisthānaparamāya)}]{\ml{അടിസ്ഥാനപരമായ} (aṭisthānaparamāya) — fundamental; basic}

\label{lesson:ML-C445-atisthanaparamaya}



\begin{quote}
Before the new one: say the Malayalam for inevitable, then the Malayalam for crucial.
\end{quote}

\subsection*{You'll want to know: \ml{അടിസ്ഥാനപരമായ}}
\textbf{\ml{അടിസ്ഥാനപരമായ}} — \emph{aṭisthānaparamāya} — \textquotedblleft{}fundamental; basic\textquotedblright{}.

\begin{quote}
\ml{അവർ} \ml{തമ്മിൽ} \ml{അടിസ്ഥാനപരമായ} \ml{ഒരു} \ml{വ്യത്യാസം} \ml{ഉണ്ട്}.
\end{quote}

\begin{quote}
\emph{avar tammil aṭisthānaparamāya oru vyatyāsaṁ uṇṭŭ.}
\end{quote}

\begin{quote}
There is a fundamental difference between them.
\end{quote}

Built on \ml{അടിസ്ഥാനം}, a basis, with -\ml{പരമായ}.

\begin{cousinweb}
\textbf{\ml{അടിസ്ഥാനപരമായ}} (\emph{aṭisthānaparamāya}) is built on \ml{അടിസ്ഥാനം}, \textquotedblleft{}a basis\textquotedblright{}, which joins native \emph{aṭi}, \textquotedblleft{}sole, bottom\textquotedblright{} (DEDR 72), to Sanskrit \emph{sthāna}, \textquotedblleft{}place\textquotedblright{} (CDIAL 13753), with -\ml{പരമായ}, \textquotedblleft{}concerned with\textquotedblright{}.
\end{cousinweb}

\begin{grammarlens}[title={What you've built}]
One more word for saying what happens if.
\end{grammarlens}

\subsection*{Your turn}
\begin{itemize}
\raggedright
  \item \emph{Say it:} \emph{aṭisthānaparamāya}
  \item \emph{Say it:} \emph{aṭisthānaparamāya}, once more
  \item \emph{Say it:} \emph{avar tammil aṭisthānaparamāya oru vyatyāsaṁ uṇṭŭ.} — There is a fundamental difference between them.
  \item \emph{Recall:} say the Malayalam for inevitable, then the Malayalam for crucial, then say \emph{aṭisthānaparamāya} again
\end{itemize}

\begin{tcolorbox}[breakable,colback=teal!4,colframe=teal!35!black,title={Before you move on}]
What does \ml{അടിസ്ഥാനപരമായ} mean? (\textquotedblleft{}Fundamental; basic\textquotedblright{}.) Say it once more.
\end{tcolorbox}
\section[\texorpdfstring{\ml{ഗുരുതരമായ} (gurutaramāya)}{ഗുരുതരമായ (gurutaramāya)}]{\ml{ഗുരുതരമായ} (gurutaramāya) — serious; grave}

\label{lesson:ML-C445-gurutaramaya}



\begin{quote}
Before the new one: say the Malayalam for crucial, then the Malayalam for fundamental.
\end{quote}

\subsection*{You'll want to know: \ml{ഗുരുതരമായ}}
\textbf{\ml{ഗുരുതരമായ}} — \emph{gurutaramāya} — \textquotedblleft{}serious; grave\textquotedblright{}.

\begin{quote}
\ml{അത്} \ml{ഗുരുതരമായ} \ml{ഒരു} \ml{തെറ്റായിരുന്നു}.
\end{quote}

\begin{quote}
\emph{atŭ gurutaramāya oru teṟṟāyirunnu.}
\end{quote}

\begin{quote}
That was a serious mistake.
\end{quote}

Said of an illness or a mistake, not of a person's mood.

\begin{cousinweb}
\textbf{\ml{ഗുരുതരമായ}} (\emph{gurutaramāya}) is built from Sanskrit \emph{gurutara}, \textquotedblleft{}heavier, very heavy\textquotedblright{} (Monier-Williams), the comparative of \emph{guru}, \textquotedblleft{}heavy\textquotedblright{}, with the adjective ending -\ml{ആയ}: a grave matter weighs heavy.
\end{cousinweb}

\begin{grammarlens}[title={What you've built}]
One more word for saying what happens if.
\end{grammarlens}

\subsection*{Your turn}
\begin{itemize}
\raggedright
  \item \emph{Say it:} \emph{gurutaramāya}
  \item \emph{Say it:} \emph{gurutaramāya}, once more
  \item \emph{Say it:} \emph{atŭ gurutaramāya oru teṟṟāyirunnu.} — That was a serious mistake.
  \item \emph{Recall:} say the Malayalam for crucial, then the Malayalam for fundamental, then say \emph{gurutaramāya} again
\end{itemize}

\begin{tcolorbox}[breakable,colback=teal!4,colframe=teal!35!black,title={Before you move on}]
What does \ml{ഗുരുതരമായ} mean? (\textquotedblleft{}Serious; grave\textquotedblright{}.) Say it once more.
\end{tcolorbox}
\section[\texorpdfstring{\ml{സങ്കീർണ്ണമായ} (saṅkīrṇṇamāya)}{സങ്കീർണ്ണമായ (saṅkīrṇṇamāya)}]{\ml{സങ്കീർണ്ണമായ} (saṅkīrṇṇamāya) — complex; complicated}

\label{lesson:ML-C445-sankirnnamaya}



\begin{quote}
Before the new one: say the Malayalam for fundamental, then the Malayalam for serious.
\end{quote}

\subsection*{You'll want to know: \ml{സങ്കീർണ്ണമായ}}
\textbf{\ml{സങ്കീർണ്ണമായ}} — \emph{saṅkīrṇṇamāya} — \textquotedblleft{}complex; complicated\textquotedblright{}.

\begin{quote}
\ml{ഇത്} \ml{വളരെ} \ml{സങ്കീർണ്ണമായ} \ml{ഒരു} \ml{പ്രശ്നമാണ്}.
\end{quote}

\begin{quote}
\emph{itŭ vaḷare saṅkīrṇṇamāya oru praśnamāṇŭ.}
\end{quote}

\begin{quote}
This is a very complex problem.
\end{quote}

The opposite of \ml{ലളിതമായ}, simple.

\begin{cousinweb}
\textbf{\ml{സങ്കീർണ്ണമായ}} (\emph{saṅkīrṇṇamāya}) is built from Sanskrit \emph{saṁkīrṇa}, \textquotedblleft{}poured together, mixed, confused\textquotedblright{} (Monier-Williams), with the adjective ending -\ml{ആയ}: a complex thing has many parts mixed.
\end{cousinweb}

\begin{grammarlens}[title={What you've built}]
One more word for saying what happens if.
\end{grammarlens}

\subsection*{Your turn}
\begin{itemize}
\raggedright
  \item \emph{Say it:} \emph{saṅkīrṇṇamāya}
  \item \emph{Say it:} \emph{saṅkīrṇṇamāya}, once more
  \item \emph{Say it:} \emph{itŭ vaḷare saṅkīrṇṇamāya oru praśnamāṇŭ.} — This is a very complex problem.
  \item \emph{Recall:} say the Malayalam for fundamental, then the Malayalam for serious, then say \emph{saṅkīrṇṇamāya} again
\end{itemize}

\begin{tcolorbox}[breakable,colback=teal!4,colframe=teal!35!black,title={Before you move on}]
What does \ml{സങ്കീർണ്ണമായ} mean? (\textquotedblleft{}Complex; complicated\textquotedblright{}.) Say it once more.
\end{tcolorbox}
